\documentclass[11pt,letterpaper]{article}
\usepackage{jahnel}
\DeclareMathOperator{\diam}{diam}
\DeclareMathOperator{\dist}{dist}

\runningtitle{The structure walk}
\hypersetup{pdftitle={The structure walk: a certificate for minimal witnesses, with five test families},pdfauthor={Christian Hirsch, Benedikt Jahnel}}

\title{\textbf{The structure walk:\\ a certificate for minimal witnesses,\\ with five test families}\\[0.6em] \large Working note (Step~1 of the plan)}

\author{%
  Christian Hirsch\thanks{Institut f\"ur Mathematische Stochastik, Georg-August-Universit\"at G\"ottingen, Goldschmidtstra\ss e~7, 37077 G\"ottingen, Germany. \texttt{christian.hirsch@uni-goettingen.de}. \orcid{0000-0003-4136-3740}.}
  \and
  Benedikt Jahnel\thanks{Institut f\"ur Mathematische Stochastik, Technische Universit\"at Braunschweig, Universit\"atsplatz~2, 38106 Braunschweig, Germany, and Weierstrass Institute for Applied Analysis and Stochastics, Mohrenstra\ss e~39, 10117 Berlin, Germany. \texttt{benedikt.jahnel@tu-braunschweig.de}. \orcid{0000-0002-4212-0065}.}%
}

\date{\today}

\begin{document}

\maketitle
\thispagestyle{plain}

\begin{abstract}
The single-cluster bound of the first-moment programme is at most the expected number of minimal witness sets of lines, a geometric integral over line arrangements. This note defines the certificate by which that integral is to be bounded: the \emph{structure walk}, which cuts the geodesic of a minimal witness into stages, each governed by the set of lines through its first box, and records only the data that is not determined by the lines already integrated. We prove the two lemmas the certificate needs, a bound on the size of the core of a structure and the measure bound with telescoping costs, formulate the local balance that has to hold at every stage, and run the five test families, bundle, star chain, relay, cross and fan, through it; all five balance. Turns of the geodesic are handled by the construction. Returns of the geodesic to an earlier structure are the one configuration for which the local balance has a factor larger than one, and we state the global argument that should absorb them. The structure lemma itself is not proved.
\end{abstract}

% =================================================================
\section{Setting}
\label{sec:setting}

Notation is as in~\cite{DEC26}: $\Lambda$ is a finite set of lines, $\Theta(\Lambda)$ the set of boxes met by at least $j_0$ lines of~$\Lambda$, $C(\Lambda)$ the $2$-connected component of~$0$ in $\Theta(\Lambda)$, and $\Lambda$ is a \emph{minimal witness} for~$d$ if $C(\Lambda)$ contains a box at sup-distance at least $d/2$ from~$0$ and no proper subset of~$\Lambda$ has this property. By~\cite[Proposition~6.2]{DEC26},
\begin{equation}
\label{eq:witness}
  \P\bigl(\diam C(0)\ge d\bigr)\le\sum_{n\ge j_0}\frac1{n!}\int\1\bigl\{\{\ell_1,\dots,\ell_n\}\text{ is a minimal witness}\bigr\}\dd\nu^{\otimes n}\,,
\end{equation}
and by~\cite[Lemma~7.1]{DEC26} every line of a minimal witness meets any fixed geodesic $\gamma=(z_0,\dots,z_m)$ of $C(\Lambda)$ from~$0$ to a box~$w$ with $\lVert w\rVert_\infty\ge d/2$ at a box with exactly $j_0$ lines. We fix $\gamma$ to be the lexicographically first such geodesic, a function of~$\Lambda$. For a box~$z$ let $S(z)=S_\Lambda(z)$ be the set of lines of~$\Lambda$ meeting $Q(z)$. We write $N(z)=\lvert S(z)\rvert$ and $\mu_1=4\mu/\pi$.

The goal is a bound on the right-hand side of~\eqref{eq:witness} of order $d^{-\beta}$. The integral is over arrangements; the difficulty is to organise it so that no sum over positions of boxes that are functions of the arrangement is performed.

% =================================================================
\section{The structure walk}
\label{sec:walk}

\begin{definition}[Structures, cores, stages]
\label{def:walk}
Let $\Lambda$ be a minimal witness with geodesic~$\gamma$. The \emph{structure} at a box~$z$ is $S(z)$, and its \emph{core} is
\begin{equation*}
  K(z)\coloneqq\{z'\in\Z^2\colon\text{at least }j_0\text{ lines of }S(z)\text{ meet }Q(z')\}\ni z\,.
\end{equation*}
The \emph{stages} of~$\Lambda$ are defined inductively. Stage~$1$ starts at $t_1\coloneqq0$. If stage~$k$ starts at the index~$t_k$, its structure is $S_k\coloneqq S(z_{t_k})$, its \emph{run} is the maximal interval $[t_k,t'_k]$ of indices such that $z_t\in K(z_{t_k})$ for all $t_k\le t\le t'_k$, and it ends at $t'_k$. If $t'_k=m$ the walk stops with $p\coloneqq k$ stages. Otherwise stage $k+1$ starts at $t_{k+1}\coloneqq t'_k+1$. The \emph{anchors} of the walk are the boxes $z_{t_k}$ and $z_{t'_k}$, $k\le p$, and the final box $z_m=w$.
\end{definition}

By construction $z_{t_{k+1}}\notin K(z_{t_k})$, so fewer than $j_0$ lines of~$S_k$ meet $z_{t_{k+1}}$, and since $N(z_{t_{k+1}})\ge j_0$ there is a line in $S_{k+1}\setminus S_k$. Such a line is \emph{fresh} at stage $k+1$ if it meets no box $z_t$ with $t<t_{k+1}$, and \emph{old} otherwise. A stage is \emph{paid} if it starts with a fresh line and \emph{unpaid} otherwise. A line is \emph{present} at an anchor if it meets that box; the \emph{span} $\sigma(\ell)$ of a line is the Euclidean distance between the first and the last anchor at which it is present (zero if there is only one), and a line \emph{leaves} at its last anchor if it is not within distance~$1$ of any other line of its last structure beyond that anchor.

\begin{lemma}[Size of a core]
\label{lem:core}
Let $S=S(z)$ with $\lvert S\rvert=N$. For every $s\ge0$ the number of boxes of $K(z)$ at sup-distance in $[s,s+1)$ from~$z$ is at most $8\binom N{j_0}$, and the boxes of~$K(z)$ at distance in $[s,s+1)$ that are met by a given $j_0$-subset $S'\subseteq S$ number at most~$8$.
\end{lemma}

\begin{proof}
All lines of~$S$ meet the unit box $Q(z)$. Two lines of $S'$ at angle $\alpha>0$ have, at distance~$s$ from~$z$, lateral separation at least $s\sin\alpha-\sqrt2$, and a box met by both has side~$1$, so a box at distance~$s$ met by all lines of~$S'$ lies within distance $1+\sqrt2$ of each of them; the set of points at distance in $[s,s+1)$ from~$z$ and within $1+\sqrt2$ of a given line is a union of two segments of length at most~$2(1+\sqrt2)\cdot\sqrt2$, which meet at most $8$ boxes. A box of $K(z)$ is met by all lines of some $j_0$-subset, and there are $\binom N{j_0}$ subsets.
\end{proof}

Thus the run of a stage with $N$ lines has, at distance~$s$ from its start, at most $8\binom N{j_0}$ candidate boxes, and the \emph{end} $z_{t'_k}$ of the run, together with the next start $z_{t_{k+1}}$ within distance~$2$ of it, is located up to a factor $24\cdot8\binom N{j_0}(1+s)$ by the distance~$s$ alone.

\begin{lemma}[Measure given the anchors]
\label{lem:measure}
Fix the anchors of the walk as boxes in the plane and, for every line, the first and the last anchor at which it is present, and whether it leaves there. The $\nu^{\otimes n}$-measure of the $n$-tuples with these data is at most
\begin{equation}
\label{eq:cost}
  \prod_{\ell\in\Lambda}c(\ell),\qquad
  c(\ell)\coloneqq\begin{cases}\mu_1,&\sigma(\ell)=0,\\[2pt] 6\mu/(1+\sigma(\ell)),&\sigma(\ell)>0,\ \ell\text{ does not leave},\\[2pt] 21\mu/(1+\sigma(\ell))^3,&\sigma(\ell)\ge4,\ \ell\text{ leaves}.\end{cases}
\end{equation}
Moreover, for every line and every sequence of anchors $u_a=u_{(0)},u_{(1)},\dots,u_{(r)}=u_b$ at which it is present, with $D_i\coloneqq\lvert u_{(i)}-u_a\rvert$,
\begin{equation}
\label{eq:telescope}
  \frac1{1+\sigma(\ell)}=\prod_{i=0}^{r-1}\frac{1+D_i}{1+D_{i+1}}\,.
\end{equation}
\end{lemma}

\begin{proof}
The three cases are the measure of lines meeting a box, the two-box bound~\cite[Lemma~3.2]{HJ26b}, and the exit cost~\cite[Lemma~7.2]{DEC26}. The identity~\eqref{eq:telescope} is a telescoping product with $D_0=0$ and $D_r=\sigma(\ell)$.
\end{proof}

The factor $(1+D_i)/(1+D_{i+1})$ is the part of the cost of~$\ell$ that is \emph{released} at the step from anchor $u_{(i)}$ to $u_{(i+1)}$. It is smaller than one when the walk moves away from the line's first anchor and larger than one when it moves back towards it.

% =================================================================
\section{The local balance}
\label{sec:balance}

\paragraph{What carries entropy.} Integrating the lines in the order of the stages at which they are fresh, the data of stage $k+1$ that are not determined by the lines integrated so far are: the distance $s_k=\lvert z_{t'_k}-z_{t_k}\rvert$ of the end of the run from its start, the position of the end among the $8\binom{N_k}{j_0}$ candidates at that distance (\cref{lem:core}), the step to the next start ($24$ choices), and the number of fresh lines there. The identity of an old line that enters at $z_{t_{k+1}}$, and the box at which it crosses the current structure, are functions of the integrated lines and carry no entropy. In particular an unpaid stage whose start is the crossing of the current core by an old line has entropy one, apart from the choice between the finitely many old lines, which is a function of the arrangement as well once the end of the run is fixed to be the first crossing after the start.

\paragraph{The balance at a stage.} Let stage~$k$ have $N_k$ lines and run length $s=s_k$, and let $F_{k+1}$ be the number of fresh lines at stage $k+1$. The entropy released at the step is at most $192\binom{N_k}{j_0}(1+s)$. The payment is the product of the fresh costs $\mu_1^{F_{k+1}}$, of the released factors~\eqref{eq:telescope} of all lines present at $z_{t'_k}$, and of the exit improvements of lines that leave at $z_{t'_k}$. The \emph{local balance} asks that, for a constant $\rho<1$ and after summing over~$s$,
\begin{equation}
\label{eq:balance}
  \sum_{s\ge0}192\binom{N_k}{j_0}(1+s)\cdot\mu_1^{F_{k+1}}\cdot\prod_{\ell\text{ present at }z_{t'_k}}\frac{1+D_\ell(t_k)}{1+D_\ell(t'_k)}\cdot\prod_{\ell\text{ leaves}}\frac{3.5}{(1+s)^2}\ \le\ \rho\,,
\end{equation}
whenever the stage is paid, and $\le C$ whenever it is unpaid, with the number of unpaid stages between two paid ones bounded by a function of the structures involved. If~\eqref{eq:balance} holds, the sum over certificates is at most $\sum_p\rho^{p'}C^{p-p'}$ over the numbers of paid and unpaid stages, and the diameter constraint, which forces either many paid stages or a run of length at least $d/(Cp)$, gives $d^{-(j_0-2)}$ up to constants, by the same computation as in~\cite[Section~3]{SC26}.

\paragraph{Why turns are not a separate case.} If the geodesic leaves a structure whose lines continue, the next box is outside the core, so the stage ends, the structure's $j_0$ lines through the end box have the step from $z_{t_k}$ to $z_{t'_k}$ released, and each of them contributes $(1+D(t_k))/(1+D(t'_k))$, which for lines fresh at $z_{t_k}$ is $1/(1+s)$. With $j_0$ such lines the sum over~$s$ in~\eqref{eq:balance} is $\sum_s192(1+s)^{1-j_0}<\infty$, and the fresh line at the next start provides the factor $\mu_1$. The case of~\cite[Section~8]{DEC26}, one line ending while the others continue, does not arise: the stage ends for all lines of the structure at once, and a line that continues with the geodesic is simply present in the next structure and releases its factor there.

% =================================================================
\section{Five test families}
\label{sec:tests}

\begin{enumerate}
  \item \emph{Bundle.} $j_0$ lines through $Q(0)$ within distance~$1$ of each other up to~$w$. One stage, $N_1=j_0$, run length $s\ge d/2$, all lines present at the final anchor~$w$ with span~$s$: cost $(6\mu/(1+s))^{j_0}$, entropy $8(1+s)$. The contribution is $\sum_{s\ge d/2}8(1+s)(6\mu)^{j_0}(1+s)^{-j_0}=O(d^{2-j_0})$, one power worse than the true order because the direction of the first line is not used to locate~$w$; see~\cite[Lemma~6.4]{DEC26}.
  \item \emph{Star chain.} Boxes $z_0,\dots,z_m$ with $j_0$ fresh lines each. Every stage has run length~$0$, $N_k=j_0$, $F_{k+1}=j_0$; entropy $24\cdot8$, payment $\mu_1^{j_0}$; contribution $(192\mu_1^{j_0})^m$. Balance~\eqref{eq:balance} holds with $\rho=192\mu_1^{j_0}$.
  \item \emph{Relay.} $j_0-1$ bundle lines through $Q(0)$ and companions $\ell'_1,\ell'_2,\dots$, companion~$i$ present on the bundle over $[x_{i-1},x_i]$ and leaving at $x_i$. Stage~$i$ starts at $x_{i-1}$ with $S_i$ the bundle lines and $\ell'_i$ (and $\ell'_{i-1}$ if it has not yet left), $N_i\le j_0+1$, and ends where $\ell'_i$ leaves, at distance $s_i=x_i-x_{i-1}$; the next stage starts with the fresh $\ell'_{i+1}$. Released at the step: $\ell'_i$ leaves, factor $3.5/(1+s_i)^2$ on top of $1/(1+s_i)$; the bundle lines release $(1+x_{i-1})/(1+x_i)$ each, which telescopes to $(1+x_k)^{-(j_0-1)}$ over the walk; $\ell'_{i+1}$ pays $\mu_1$. Balance: $\sum_s192(j_0+1)(1+s)\cdot3.5(1+s)^{-3}\mu_1\le\rho$ for $\mu_1$ small, and the total is $O(d^{-(j_0-2)})$, in agreement with~\cite[Lemma~7.3]{DEC26}. Without the exit cost the sum over~$s$ would be $\sum_s(1+s)\cdot(1+s)^{-1}$, divergent: the exit cost is what makes relays balance.
  \item \emph{Cross.} Bundle~$B_1$ through $Q(0)$, bundle~$B_2$ through a box~$x$ of~$B_1$ at an angle, the geodesic riding $B_1$ to~$x$ and $B_2$ to~$w$. Stage~$1$: $S_1=B_1$, run to~$x$ (the next geodesic box lies on $B_2$ off the core of~$B_1$), $s_1=\lvert x\rvert$, released $(1+s_1)^{-j_0}$, entropy $8(1+s_1)$; stage~$2$ starts at the first box of~$B_2$ after~$x$ with $j_0$ fresh lines, $\mu_1^{j_0}$, run to~$w$, released $(1+s_2)^{-j_0}$ at the final anchor, entropy $8(1+s_2)$. Contribution $\sum_{s_1+s_2\ge d/2}64(1+s_1)^{1-j_0}(1+s_2)^{1-j_0}\mu_1^{j_0}(6\mu)^{2j_0}=O(d^{2-j_0})$.
  \item \emph{Fan.} $N\ge j_0+1$ lines through $Q(0)$ with distinct small angles, the geodesic riding the core, in which a sliding window of $j_0$ lines is present at distance~$s$; all $N$ lines are critical somewhere, so the fan is a minimal witness. One stage, $N_1=N$, fresh cost $\mu_1^N$, entropy $8\binom N{j_0}(1+s)$, released at~$w$: the $j_0$ lines present there have span~$s$, the others have shorter spans $\sigma_\ell$ with $(6\mu/(1+\sigma_\ell))\le\mu_1$. Contribution at most $\sum_s8\binom N{j_0}(1+s)(6\mu)^{j_0}(1+s)^{-j_0}\mu_1^{N-j_0}$, and $\binom N{j_0}\mu_1^{N-j_0}\le(2\mu_1)^{N-j_0}2^{j_0}$ is summable over~$N$ for $\mu_1<1/2$. The fan balances.
\end{enumerate}

% =================================================================
\section{Returns}
\label{sec:returns}

The released factor~\eqref{eq:telescope} is larger than one when the walk moves back towards a line's first anchor. This happens when the geodesic returns to a structure it has left: a line~$\ell$ of that structure is present at an early anchor $u_a$, absent during an excursion, and present again at a later anchor~$u_b$ closer to~$u_a$ than the anchors of the excursion. Globally the product over the steps is still $1/(1+\sigma(\ell))\le1$, so the measure bound is correct; what fails is the local form~\eqref{eq:balance}, which cannot be summed step by step when some steps have factors above one.

The proposed treatment is to group stages into \emph{excursions}: a maximal sequence of stages during which the walk is outside the $2$-neighbourhood of the cores of all earlier structures. Within an excursion all released factors of the lines of earlier structures are at most one, since the walk moves away from the structures it left, up to a bounded factor accounting for the non-monotonicity of Euclidean distance along a $2$-path that stays outside a neighbourhood; the stages of an excursion are balanced by~\eqref{eq:balance}. A return is the first box after an excursion that lies in the core of an earlier structure. The return box is determined by the earlier structure and by the excursion's last line, which crosses that core, so it carries no entropy; and the lines of the earlier structure release at the return the factor $(1+D(\text{exit}))/(1+D(\text{return}))$, which is at most the ratio of the distances of the exit and the return from the structure's start, a factor bounded by the excursion's own displacement, and therefore paid by the excursion if the excursion's balance is strict enough, which requires the fresh lines of the excursion to pay for its displacement. This is where the argument currently stops: an excursion that makes a large displacement with few fresh lines is a chain of long bundles, whose own balance pays only for their lengths, not for an extra factor equal to the displacement. We expect this to be resolved by noting that a return to an earlier structure means that the component $C(\Lambda)$ contains a loop, so that the lines of the excursion, or those of the earlier structure beyond the exit, are not all critical, contradicting minimality; a precise version of this is the statement that the structure walk of a minimal witness has no returns at all, which would remove the case. We have not proved it.

% =================================================================
\section{Step 2: returns}
\label{sec:step2}

We now remove the case of \cref{sec:returns}. The first observation is that the factors larger than one are produced by the definition of the span, not by the geometry.

\begin{definition}[Running maximum]
\label{def:runmax}
For a line~$\ell$ and an anchor~$u$ at which $\ell$ is present, let $M_\ell(u)$ be the largest Euclidean distance between two anchors at which~$\ell$ is present, among the anchors up to and including~$u$ in walk order. The \emph{span} of~$\ell$ is $\sigma(\ell)\coloneqq M_\ell(u_{\mathrm{last}})$, and the factor \emph{released} by~$\ell$ at a presence~$u$ following its previous presence~$u^-$ is $(1+M_\ell(u^-))/(1+M_\ell(u))$.
\end{definition}

\begin{lemma}[Monotone release]
\label{lem:monotone}
With \cref{def:runmax}, every released factor is at most one, the product of the released factors of~$\ell$ over its presences is $1/(1+\sigma(\ell))$, and \cref{lem:measure} holds with this span, with the same costs~\eqref{eq:cost}.
\end{lemma}

\begin{proof}
$M_\ell$ is non-decreasing in walk order, so each factor is at most one, and the product telescopes from $M_\ell=0$ at the first presence to $\sigma(\ell)$. The line passes through the two anchors realising $\sigma(\ell)$, which are among the fixed anchor data, so the two-box bound and the exit cost apply with $\sigma(\ell)$.
\end{proof}

A return to an old line therefore costs nothing: the walk arrives at the return box by the stage it is in, whose entropy is counted there, and the presence of the old line at that box is a constraint that releases a factor at most one. What a return \emph{can} do is create a stage served only by old lines, with no fresh line to pay for it, and we now show that the canonical geodesic can be chosen so that this happens only at boxes determined by the arrangement.

\paragraph{The widened core and the canonical geodesic.} Replace $K(z)$ in \cref{def:walk} by the set $\widetilde K(z)$ of boxes at sup-distance at most~$1$ from a box met by at least $j_0$ lines of $S(z)$. \Cref{lem:core} holds for $\widetilde K$ with $8$ replaced by~$72$, and a box outside $\widetilde K(z)$ is still met by fewer than $j_0$ lines of $S(z)$, so the stages are defined as before. Fix~$w$ as the lexicographically first box of $C(\Lambda)$ at sup-distance at least $d/2$ from~$0$. Call a box~$y$ \emph{admissible} from a box~$x$ of~$\gamma$ if $y\in\Theta(\Lambda)$, $\lVert y-x\rVert_\infty\le2$ and $\dist_\Theta(y,w)=\dist_\Theta(x,w)-1$, where $\dist_\Theta$ is the graph distance in the $2$-adjacency graph of $\Theta(\Lambda)$. The \emph{canonical geodesic} is built step by step: from the current box, with current structure~$S_k$, go to the lexicographically first admissible box in $\widetilde K(z_{t_k})$ if there is one, and otherwise to the lexicographically first admissible box. This is a shortest path from~$0$ to~$w$, it is a function of~$\Lambda$, and the stages are read off from it as in \cref{def:walk}.

\begin{lemma}[Tubes]
\label{lem:tube}
Let $\ell_1,\dots,\ell_{j_0}$ be lines that pass within distance~$1$ of a segment~$I$ of a line~$\ell_0$ along the whole of~$I$, and let $T$ be the set of boxes within sup-distance~$1$ of a box met by all of them. Then for any two boxes $x,y\in T$ whose projections on~$\ell_0$ lie in~$I$ and are at least $4$ apart, there is a $2$-path from~$x$ to~$y$ inside~$T$ of length $\lceil\lVert y-x\rVert_\infty/2\rceil$, and for every box~$v$ with $\lVert v-x\rVert_\infty\ge4$, the distances $\lceil\lVert v-x'\rVert_\infty/2\rceil$ are equal for all boxes $x'\in T$ with the same projection on~$\ell_0$ as~$x$.
\end{lemma}

\begin{proof}
The segment from the centre of~$x$ to the centre of~$y$ lies within distance $1+\sqrt2$ of~$\ell_0$ and within distance $2+\sqrt2$ of each~$\ell_i$. Let $v=y-x$ and assume $\lvert v_1\rvert\ge\lvert v_2\rvert$, so $\lVert v\rVert_\infty=\lvert v_1\rvert$. The boxes $x+(2k,\lfloor2kv_2/v_1\rceil)$, $0\le k\le\lvert v_1\rvert/2$, with the last step adjusted, form a $2$-path of length $\lceil\lvert v_1\rvert/2\rceil$ whose boxes are within distance~$1$ of the segment, hence within distance $3+\sqrt2<4.5$ of each~$\ell_i$; a box within distance~$4.5$ of $j_0$ lines that are within distance~$1$ of each other is within sup-distance~$1$ of a box met by all of them only if the lines pass within $2+\sqrt2$ of its centre, which holds here. For the second claim, boxes with the same projection on~$\ell_0$ differ by at most~$3$ in the direction normal to~$\ell_0$, and for $\lVert v-x\rVert_\infty\ge4$ the sup-distance to~$v$ is attained in the direction of~$\ell_0$ up to an error that does not change the value of $\lceil\cdot/2\rceil$ by more than the rounding, which we absorb by widening the tube by one more box.
\end{proof}

The lemma is stated with the constants we could verify in the time available; the second claim needs the widening mentioned at its end, which costs a factor~$9$ in \cref{lem:core} and nothing else. Its consequence is the following.

\begin{proposition}[No re-entries into the same core]
\label{prop:noreturn}
With the canonical geodesic and the widened core, let $x=z_{t'_k}$ be the end of stage~$k$. There is no index $t>t'_k+1$ such that $y=z_t$ lies in $\widetilde K(z_{t_k})$ in the tube of a $j_0$-subset of~$S_k$ that is present at both~$x$ and~$y$.
\end{proposition}

\begin{proof}
Suppose~$y$ exists. By \cref{lem:tube} there is a $2$-path inside the tube from~$x$ to~$y$ of length $\lceil\lVert y-x\rVert_\infty/2\rceil\le\dist_\Theta(x,y)$; all its boxes are met by the $j_0$ lines of the subset, so it lies in $\Theta(\Lambda)$ and is a shortest path from~$x$ to~$y$ in~$\Theta$. Concatenated with the part of~$\gamma$ from~$y$ to~$w$ it is a shortest path from~$x$ to~$w$, so its first box is admissible from~$x$ and lies in $\widetilde K(z_{t_k})$. The canonical rule then continues stage~$k$ at~$x$, contradicting that stage~$k$ ends at~$x$.
\end{proof}

\paragraph{What a stage boundary now is.} With \cref{prop:noreturn} and the rule, a stage ends at~$x$ in exactly two situations. Either some $j_0$-subset of~$S_k$ present at~$x$ stops being present within the next step, which means that a line of~$S_k$ leaves the $1$-neighbourhood of the others (exit cost), or has reached the end of its presence interval; or every admissible continuation leaves the core while the lines continue, in which case the next box carries a line not in~$S_k$ that is fresh, paying~$\mu_1$, or old. An old line entering at a box where the lines of~$S_k$ continue crosses the core of~$S_k$, and the crossing box is a function of the two lines; such a boundary carries no entropy. The second claim of \cref{lem:tube} ensures that a change of row inside the tube is never forced by admissibility except within $4$ boxes of the stage end, so lateral switches between companions, the case that produced the logarithmic loss in~\cite[Section~8]{DEC26}, do not occur before the actual end of the stage.

\paragraph{The remaining accounting.} Stages served only by old lines, produced by zero-entropy crossings and by re-entries into a \emph{different} tube of an earlier structure, inherit from \cref{lem:monotone} a released product that telescopes exactly over the \emph{super-stage}, the maximal run of consecutive stages served by the same $j_0$-subset. The local balance~\eqref{eq:balance} has to be verified per super-stage rather than per stage: its entropy is at most $\prod_kC(1+s_k/D_k)$ with $D_k\ge\sum_{i<k}s_i$ by the additivity of the running maximum along a straight tube, which is at most $C^{\#\text{stages}}(\sum_ks_k)/s_1$, polynomial in the total length, while the released product is $(1+\sum_ks_k)^{-j_0}$ up to constants. The number of zero-entropy boundaries in a super-stage is at most the number of lines crossing its tube, each of which pays~$\mu_1$ once, so $C^{\#\text{stages}}$ is absorbed for $C\mu_1<1$. We have checked this accounting on the five families and on a bundle crossed by $q$ transversal lines; a full proof of the super-stage balance, in particular of the additivity claim $D_k\ge\sum_{i<k}s_i$ for re-entries into different tubes of the same structure, is what remains for Step~3.

% =================================================================
\section{Step 3: the summation and the local balance}
\label{sec:step3}

\paragraph{Free boundaries.} By \cref{sec:step2} a stage ends at $b_k$, the last box of its run met by at least $j_0$ lines of~$S_k$, for one of two reasons: a serving line (present at $a_k$ and $b_k$) is absent from every box within distance~$1$ of the next geodesic box, which for a straight line means that it leaves the $1.5$-neighbourhood of the other serving lines within distance~$4$ of~$b_k$ (an \emph{exit}); or all serving lines continue but every admissible continuation leaves the core (a \emph{turn}). At a turn the next box carries a line not in~$S_k$. If that line is fresh, the turn is \emph{free}; if all lines of $S_{k+1}$ are old, the next box lies on the intersection of the current tube with the tube of an old structure, which is a function of the arrangement, and the turn is \emph{determined}. Exits are free. The entropy of the certificate is carried by the free boundaries only: at a determined boundary the position $b_k$ is a function of the integrated lines.

\paragraph{The per-stage sum.} Let stage~$k$ have $N_k$ lines through~$a_k$, serving lines $\mathcal S_k\subseteq S_k$ (present at $a_k$ and~$b_k$, $\lvert\mathcal S_k\rvert\ge j_0$), of which $f_k$ have no presence before~$a_k$ and $o_k$ are old, with $D_k\coloneqq\min\{M_\ell(a_k)\colon\ell\in\mathcal S_k\text{ old}\}$ the smallest running maximum among the old serving lines, $D_k\coloneqq0$ if there are none. Let $F_{k+1}$ be the number of fresh lines at $a_{k+1}$, and let $x_k=1$ if the boundary is an exit of a line fresh at~$a_k$ and $x_k=0$ otherwise. Define
\begin{equation}
\label{eq:phi}
  \Phi_k\coloneqq\sum_{s\ge0}E_k(s)(1+s)^{-f_k-x_k}\Bigl(\frac{1+D_k}{1+\max\{D_k,s\}}\Bigr)^{o_k}\mu_1^{F_{k+1}},\quad
  E_k(s)\coloneqq\begin{cases}C_0N_k(1+s),&o_k=0,\\ C_0N_k(1+s/D_k),&o_k\ge1,\end{cases}
\end{equation}
where $E_k(s)$ bounds the number of candidates for~$b_k$ at distance~$s$ from~$a_k$ and the step to~$a_{k+1}$: by \cref{lem:core} and its refinement (boxes at distance~$s$ met by $j_0$ lines through~$a_k$ are met by lines within angle $2/s$ of each other, so there are at most $16N_k$ of them per unit distance) when all serving lines are new, and by the direction of an old serving line, known to precision $1/D_k$, when there is one. The factors are the released costs of \cref{lem:monotone}: $1/(1+s)$ for a line with its first presence at~$a_k$, the extra $1/(1+s)$ of the exit cost $21\mu/(1+s)^2$ for a fresh leaver (the version of \cite[Lemma~7.2]{DEC26} with entry at~$a_k$ and exit at~$b_k$), and $(1+D_k)/(1+\max\{D_k,s\})$ for an old serving line, whose running maximum rises to at least~$s$.

\begin{theorem}[Summation]
\label{thm:sum}
Suppose that for every minimal witness and every free stage~$k$ of its walk, $\Phi_k\le\rho<1$ if $F_{k+1}\ge1$ or $x_k=1$, and that the number of determined boundaries between two consecutive free ones is at most the number of lines crossing the tube of the stage, each of which has its own cost~$\mu_1$. Then there are $C$ and $\beta\ge j_0-2$ with
\begin{equation}
\label{eq:final}
  \sum_{n}\frac{M_n}{n!}\le C\,d^{-\beta}\qquad(d\ge1)\,.
\end{equation}
\end{theorem}

\begin{proof}
Integrate the lines stage by stage. At stage~$k$ the fresh lines at~$a_k$ are integrated over the lines meeting $Q(a_k)$, measure $\mu_1$ each, which is the factor $\mu_1^{F_k}$ attributed to the boundary $k-1$; the end $b_k$ and the start $a_{k+1}$ are summed over their candidates, at most $E_k(s)$ for each~$s$, and the released factors of \cref{lem:monotone} are applied to the serving lines, whose product over all stages is at most $\prod_\ell c(\ell)/\mu_1$ by \cref{lem:measure}. Determined boundaries contribute a factor~$1$ in entropy and a released factor at most~$1$. Hence the sum over certificates with $p$ free stages is at most $\prod_k\Phi_k\le\rho^{p}$ times the contribution of the determined boundaries, which is at most $\prod(C\mu_1)^{\#\text{crossing lines}}\le1$ for $C\mu_1\le1$. The diameter constraint $\lVert w\rVert_\infty\ge d/2$ gives $\sum_ks_k\ge d/2-4p$, so either $p\ge d/16$, contributing at most $\sum_{p\ge d/16}\rho^p\le C\rho^{d/16}$, or some stage has $s_k\ge d/(4p)$ with $p<d/16$; for that stage the summand of~\eqref{eq:phi} restricted to $s\ge d/(4p)$ is at most $C(1+s)^{1-j_0}$ summed over $s\ge d/(4p)$, i.e.\ $C(d/4p)^{2-j_0}$, since $f_k+o_k\ge j_0$ and each released factor is at most $(1+D_k)/(1+s)\le(1+s)^{-1}$ once $s\ge d/(4p)\ge D_k$, the running maximum of an old line being bounded by the displacement of the walk so far, which is at most $4\sum_{i<k}s_i+4k\le4d/(4p)\cdot p$; this bound is crude but suffices. Summing $p\,\rho^{p-1}C(d/4p)^{2-j_0}$ over $p<d/16$ gives $O(d^{2-j_0})$.
\end{proof}

The proof is routine once the hypothesis is granted; the hypothesis is the content.

\begin{lemma}[Local balance, proved cases]
\label{lem:balance}
Let $\mu_1\le\mu^*$ with $\mu^*$ small enough. Then $\Phi_k\le\rho<1$ holds in each of the following cases.
\begin{enumerate}
  \item $o_k=0$ (all serving lines new) and $F_{k+1}\ge1$ or $x_k=1$: $\Phi_k\le C_0N_k\mu_1^{F_{k+1}}\sum_s(1+s)^{1-j_0-x_k}\le\rho$, since $j_0\ge4$ and $N_k\le F_k+N^{\mathrm{old}}_k$ with the old lines through~$a_k$ determined.
  \item $o_k\ge1$ and $f_k+x_k\ge2$: splitting the sum at $s=D_k$,
  \begin{equation*}
    \sum_{s\le D_k}\bigl(1+\tfrac{s}{D_k}\bigr)(1+s)^{-2}+\sum_{s>D_k}\tfrac{s}{D_k}(1+s)^{-f_k-x_k}\Bigl(\tfrac{1+D_k}{1+s}\Bigr)^{o_k}\le4+C\,D_k^{\,o_k-1}\sum_{s>D_k}s^{1-f_k-x_k-o_k}\le C'\,,
  \end{equation*}
  using $f_k+x_k+o_k\ge j_0+1\ge3$; so $\Phi_k\le C_0N_kC'\mu_1^{F_{k+1}}\le\rho$ when $F_{k+1}\ge1$, and when $F_{k+1}=0$ the boundary is an exit of a fresh line ($x_k=1$) and the serving set at the next stage has lost a line, so the next stage has a fresh line or is determined.
  \item $o_k\ge1$, $f_k+x_k\le1$, and the run exceeds the known span, $s>D_k$: the summand is at most $\frac{s}{D_k}(1+s)^{-f_k-x_k}\bigl(\frac{1+D_k}{1+s}\bigr)^{o_k}$ with $o_k\ge j_0-1$, whose sum over $s>D_k$ is at most $C\,D_k^{\,o_k-1}D_k^{\,2-f_k-x_k-o_k}=C\,D_k^{\,1-f_k-x_k}\le C$.
\end{enumerate}
In the relay, case~2 applies at every exchange with $f_k=x_k=1$: the sum over $s\le D_k$ is bounded and the sum over $s>D_k$ is paid by the release of the old bundle lines; this is where the earlier logarithmic loss disappears.
\end{lemma}

\begin{proof}
The computations are those displayed; in case~1 the sum $\sum_s(1+s)^{1-j_0-x_k}$ converges for $j_0\ge3$, and in cases 2 and 3 the bounds follow from $\sum_{s>D}s^{-a}\le CD^{1-a}$ for $a\ge2$. The claim about the next stage in case~2 follows from the definition of a stage boundary: the serving set of stage $k+1$ cannot be the serving set of stage~$k$.
\end{proof}

\paragraph{The remaining case.} Not covered is an all-old stage, $f_k+x_k\le1$ with $o_k\ge j_0-1$, whose run stays within the known span, $s\le D_k$. Its summand is $C_0N_k(1+s/D_k)(1+s)^{-f_k-x_k}$ with no release, and the sum over $s\le D_k$ is of order $D_k^{\,1-f_k-x_k}$, up to $D_k$ when $f_k=x_k=0$. Such a stage is served by at least $j_0-1$ lines each of which has two earlier presences at distance at least~$D_k$, and which now run together in one tube for length~$s\le D_k$. If these lines were jointly present at an earlier anchor, the stage is a re-entry into the core of that anchor's structure along their common tube, excluded by \cref{prop:noreturn}. Otherwise the lines come from different earlier structures and are nearly parallel, within angle $2/s$, which is a constraint of strength about $1/s$ on each of them that the pair bound through their own two anchors does not record: the earlier anchors of such a line, whose entropy was counted as $C(1+D)$ at the time, are confined by the parallelism to a sector of angular width $2/s$, so that the earlier entropy should have been $C(1+D/s)$, a refund of order $1/s$ per line, which pays the present sum. Making this refund rigorous requires integrating the lines in an order in which the parallelism is visible when the earlier anchors are summed, that is, a global rearrangement of the certificate rather than a local balance. We regard this as the last open point of the structure lemma. On the five test families and on their combinations we have not found a configuration of this type with more than a bounded number of such stages, and we conjecture that minimality excludes it: a line with two far presences that later joins a tube of lines from other structures is critical at its two earlier anchors and at the tube, three places, and the geodesic between them could be shortened along the line unless the line is thin in between, which bounds the number of such lines by the number of thin stretches the geodesic avoids.

% =================================================================
\section{Step 4: the global accounting and the conditional theorem}
\label{sec:step4}

The remaining case of \cref{sec:step3} splits into two according to how the stage ends.

\begin{lemma}[Exits of old lines carry no entropy]
\label{lem:oldexit}
Let stage~$k$ be served by old lines and end because a serving line~$\ell$ with running maximum $D=M_\ell(a_k)\ge1$ leaves the $1.5$-neighbourhood of the reference line~$\ell_0$ of the stage within distance~$4$ of~$b_k$. Then the sum over the candidates for~$b_k$ of the conditional measure of~$\ell$, given its two extreme earlier anchors and given~$\ell_0$, is at most the pair bound $6\mu/(1+D)$ times an absolute constant. Consequently such a boundary contributes a factor at most~$C$ to the certificate sum, with no sum over~$s$.
\end{lemma}

\begin{proof}
The line~$\ell$ passes through its two extreme anchors $u,u'$ at distance $D'\ge D$; parametrise it by its angle~$\alpha$ to~$\ell_0$ and its crossing~$c$ with~$\ell_0$. The exit position is $c+1.5/\sin\alpha$, a function of $(\alpha,c)$. The set of $(\alpha,c)$ compatible with $u,u'$ has $\nu$-measure at most $6\mu/(1+D')$, and the boxes $b_k$ are the unit intervals of~$\ell_0$ into which the exit falls; summing over the boxes the measure of the lines whose exit falls into each box gives the measure of all compatible lines, which is at most $6\mu/(1+D')$. The constant accounts for the step from $b_k$ to~$a_{k+1}$ and for the choice of~$\ell$ among the serving lines, at most $N_k$, which is determined by the arrangement once the exit box is fixed to be the first exit after~$a_k$.
\end{proof}

The lemma replaces the entropy count by the measure itself, which is the correct way to treat any boundary whose position is a function of lines already integrated: the exit of a known line, the crossing of two known tubes, and the junction of a known tube with a known structure are all of this kind. What is left is the one boundary whose position is not a function of the integrated lines.

\begin{definition}[Bad stages]
\label{def:bad}
A stage~$k$ is \emph{bad} if its serving lines are all old, no serving line leaves at~$b_k$, and the next box $a_{k+1}$ carries at least one fresh line: the walk rides a tube of lines it has seen before and turns off it, at a position that is free along the tube, onto a structure with fresh lines.
\end{definition}

For a bad stage, the sum over~$s$ of the released factors is $\sum_{s\le D_k}C_0N_k(1+s/D_k)\mu_1^{F_{k+1}}\le2C_0N_kD_k\mu_1^{F_{k+1}}$, and nothing else is released, so a bad stage costs up to a factor $D_k$ that has to be paid from elsewhere. The payment exists: each old serving line~$\ell$ has $\sigma(\ell)\ge D_k$, and its cost $6\mu/(1+\sigma(\ell))$ was released at its own stages against entropy that was only linear in the run length while $j_0$ lines released $(1+s)^{-j_0}$; the surplus $(1+s)^{1-j_0}$ of a fresh bundle stage is unused. The following hypothesis says that this surplus is not asked to pay too often.

\begin{definition}[Hypothesis $\mathrm{(B)}$]
\label{def:B}
Hypothesis $\mathrm{(B)}_B$ holds if, for every minimal witness, every line serves at most~$B$ bad stages of the structure walk.
\end{definition}

\begin{theorem}[Conditional structure lemma]
\label{thm:conditional}
Assume $\mathrm{(B)}_B$ with $B\le j_0-7$, let $\mu_1\le\mu^*$, and assume the tube lemma (\cref{lem:tube}) with the constants stated there. Then
\begin{equation}
\label{eq:condfinal}
  \P\bigl(\diam C(0)\ge d\bigr)\le\sum_n\frac{M_n}{n!}\le C\,d^{-(j_0-B-2)}\qquad(d\ge1)\,,
\end{equation}
so the single-cluster bound~$\mathrm{(S)}$ of~\cite{FM26} holds with $\beta=j_0-B-2>4$.
\end{theorem}

\begin{proof}
Run the summation of \cref{thm:sum} with the following modification of the payment rule. For a bad stage~$k$ assign its factor $2C_0N_kD_k\mu_1^{F_{k+1}}$ to its serving lines, each of which receives $(2C_0N_kD_k)^{1/(j_0-1)}$, which is at most $(2C_0N_k)^{1/(j_0-1)}(1+\sigma(\ell))^{1/(j_0-1)}$ since $\sigma(\ell)\ge D_k$. By $\mathrm{(B)}_B$ a line receives at most~$B$ such assignments, so its total cost after assignments is at most $C^B\,6\mu\,(1+\sigma(\ell))^{-1+B/(j_0-1)}$. All other free stages have $\Phi_k\le\rho$ by \cref{lem:balance}, all determined boundaries and exits of old lines contribute at most~$C$ by \cref{lem:oldexit} and the paragraph following it, and the number of determined boundaries is at most the number of lines crossing the tubes, each paying~$\mu_1$. The diameter constraint forces either at least $d/16$ free stages, contributing $\rho^{d/16}$, or a serving set of $j_0$ lines with $\sigma\ge d/(4p)$, contributing at most $C(d/4p)^{1-j_0+Bj_0/(j_0-1)}\cdot(d/4p)$, where the last factor is the entropy of that stage; summing over $p$ as in \cref{thm:sum} gives $d^{-(j_0-B-2)}$ up to constants, using $Bj_0/(j_0-1)\le B+1$.
\end{proof}

\paragraph{On hypothesis $\mathrm{(B)}$.} A bad stage served by~$\ell$ requires $j_0-2$ further old lines, from structures in which $\ell$ was not present, to run together with~$\ell$ in one tube; they are then within angle $2/s$ of~$\ell$, which is a constraint of strength about $1/s$ on the direction of each, not recorded by their pair bounds. After the bad stage the walk leaves the tube and, by \cref{prop:noreturn}, does not return to it; a second bad stage served by~$\ell$ needs $j_0-2$ \emph{new} old lines converging onto~$\ell$ further along. Two lines within angle $2/s$ of~$\ell$ that meet it at positions $x$ and~$x'$ with $\lvert x-x'\rvert\gg s$ cross each other near~$\ell$, so the converging lines of distinct bad stages on~$\ell$ are pairwise nearly parallel and pairwise crossing near~$\ell$: a family of $(j_0-2)B$ lines all within angle $2/s$ of a common line and all passing within distance~$1$ of it, which is a bundle of width~$1$ and length~$Bs$ containing~$\ell$. Such a bundle makes the stretch of~$\ell$ between the bad stages thick, and the canonical geodesic, which prefers the current core and whose shortest in-tube paths exist by \cref{lem:tube}, would not have left~$\ell$ between them. We therefore expect $\mathrm{(B)}_B$ to hold with $B=1$, but the argument as sketched needs the converging lines to be within distance~$1$ of~$\ell$ along the whole stretch, which the angle bound gives only up to distance $Bs\cdot2/s=2B$, and closing this requires a quantitative version with the tube width depending on~$B$. We have not completed it.

\paragraph{Status of the whole chain.} Granting $\mathrm{(B)}_B$, \cref{thm:conditional} gives the single-cluster bound. The first-moment note~\cite[Section~3.1]{FM26} derives the thick-cluster lemma~$\mathrm{(T)}$ from~$\mathrm{(S)}$ by a count that treats distinct clusters as independent and says so; that step has to be carried out with the Mecke formula on the witnessing lines of a chain of clusters, as outlined there, and is a second open link. The conditional theorem~\cite[Theorem~1.2]{FM26} then gives $\lambda_c(\mu,1)>0$. So the subcritical phase in the sparse regime currently rests on two statements: $\mathrm{(B)}$ above, a geometric property of minimal arrangements, and the dependence step $\mathrm{(S)}\Rightarrow\mathrm{(T)}$. Neither is proved.

% =================================================================
\section{Summary}
\label{sec:summary}

\begin{enumerate}
  \item The structure walk (\cref{def:walk}) cuts the geodesic of a minimal witness into stages governed by the set of lines through the stage's first box; its anchors are the starts and ends of the runs and the final box. \Cref{lem:core,lem:measure} are rigorous: the core of a structure with $N$ lines has at most $8\binom N{j_0}$ boxes at each distance, and the measure of the arrangements with given anchor data is at most the product of the per-line costs~\eqref{eq:cost}, which distribute over the stages by~\eqref{eq:telescope}.
  \item The local balance~\eqref{eq:balance} is the statement to be proved. It holds on the five test families, bundle, star chain, relay, cross and fan, and the construction makes turns of the geodesic a non-case. The relay requires the exit cost; the fan requires the sum over the number of lines through a box, which is paid by the fresh costs.
  \item \Cref{sec:step2} (Step~2) removes the return case: with the span defined through a running maximum, all released factors are at most one (\cref{lem:monotone}); with a canonical geodesic that prefers the current widened core, the walk never re-enters a core along a tube it has left (\cref{prop:noreturn}), and lateral switches inside a tube are excluded by the row-independence of distances (\cref{lem:tube}). Stage boundaries are then exits, fresh lines, or zero-entropy crossings by old lines. \Cref{sec:step3} (Step~3) proves the summation over certificates conditionally on a local balance at the free stage boundaries (\cref{thm:sum}) and proves the balance in all cases except one (\cref{lem:balance}): an all-old stage whose run stays within the known span of its lines, with the lines coming from different earlier structures. \Cref{sec:step4} (Step~4) shows that exits of old lines carry no entropy (\cref{lem:oldexit}), isolates the \emph{bad stages} (free turns out of tubes of old lines), and proves the structure lemma conditionally on hypothesis $\mathrm{(B)}_B$, that no line serves more than~$B$ bad stages (\cref{thm:conditional}, exponent $j_0-B-2$). We give the geometric reason to expect $\mathrm{(B)}_1$ and the quantitative gap in it. The chain to the subcritical phase then has exactly two unproved links: $\mathrm{(B)}$, and the dependence step from~$\mathrm{(S)}$ to the thick-cluster lemma~$\mathrm{(T)}$ in~\cite{FM26}.
\end{enumerate}

\section*{Acknowledgements}

BJ acknowledges the financial support of the Leibniz Association within the Leibniz Junior Research Group on Probabilistic Methods for Dynamic Communication Networks (DYCOMNET) as part of the Leibniz Competition (grant no.~J105/2020), of the Berlin Mathematics Research Center MATH+ through the project EF-MA-Sys-2 on Information Flow \& Emergent Behavior in Complex Networks, and of the Deutsche Forschungsgemeinschaft (DFG) through Project P27 within the SPP~2265.

\section*{Use of AI tools}

The simulations, proofs, computations, figures and the text of this note were developed with substantial assistance of the AI systems Claude (Anthropic) and ChatGPT Astra~6 (OpenAI), in an interactive process directed by the authors.

\bibliographystyle{amsalpha}
\bibliography{references}

\end{document}
